
\section{Experimental Setup}
Experiments were run on a Windows development machine using the scripts in \texttt{run\_evaluation\_pipeline.py}. Models: merged/LoRA Bloom classifier (\texttt{models/qwen\_bloom\_merged} or \texttt{qwen\_bloom\_3000}), Qwen2.5-1.5B-Instruct Q4\_K\_M GGUF ($\approx$986\,MB), all-MiniLM-L6-v2 embeddings, FAISS CPU index. Offline flags \texttt{HF\_HUB\_OFFLINE=1} were used where applicable.

\section{Performance Metrics}
\begin{itemize}
    \item \textbf{Bloom:} accuracy, macro-F1, weighted-F1, mean ordinal distance, within-one-level accuracy, severe-error rate.
    \item \textbf{RAG:} exact match, token-F1, retrieval hit@1, hit@3, unsupported-answer proxy.
    \item \textbf{Privacy:} attack block rate, benign allow rate, per-category block rate.
    \item \textbf{Multimodal:} answer contains expected substring (smoke tests).
\end{itemize}

\section{Result Analysis}
\subsection{Bloom Taxonomy Classification}
Table~\ref{tab:bloom_baseline_comparison} consolidates metrics from \texttt{evaluation\_outputs/} (TF-IDF + LinearSVC and zero-shot Qwen GGUF on a 15\% stratified hold-out, $n{\approx}379$) and \texttt{evaluation\_results/metrics.json} (fine-tuned LoRA on the official Figshare test split, $n{=}2330$). Artifacts are merged by \texttt{build\_bloom\_comparison.py} into \texttt{results/bloom\_baseline\_comparison.json}.

\begin{table}[H]
\centering
\caption{Bloom baseline comparison (\texttt{build\_bloom\_comparison.py})}
\label{tab:bloom_baseline_comparison}
\renewcommand{\arraystretch}{1.2}
\begin{tabular}{|l|c|c|c|}
\hline
\textbf{Metric} & \textbf{SVM (hold-out)} & \textbf{Zero-shot GGUF (hold-out)} & \textbf{LoRA (official test)} \\
\hline
Accuracy & 0.839 & 0.441 & 0.748 \\
\hline
Macro-F1 & 0.826 & 0.369 & 0.721 \\
\hline
Within-one-level acc. & 0.916 & 0.686 & 0.880 \\
\hline
Severe-error rate & 0.084 & 0.314 & 0.064 \\
\hline
Mean ordinal distance & 0.322 & 1.074 & 0.452 \\
\hline
\end{tabular}
\end{table}

\noindent SVM vs.\ zero-shot label agreement on the hold-out split was 47.5\% (\texttt{evaluation\_outputs/evaluation\_report.json}). The trained LoRA model outperforms zero-shot prompting and approaches the lexical SVM on accuracy while generalizing on the larger official test set.

\begin{figure}[H]
\centering
\begin{subfigure}{0.48\textwidth}
    \centering
    \includegraphics[width=\linewidth]{images/bloom_svm_cm.png}
    \caption{SVM baseline}
    \label{fig:bloom_svm_cm}
\end{subfigure}
\hfill
\begin{subfigure}{0.48\textwidth}
    \centering
    \includegraphics[width=\linewidth]{images/bloom_zero_shot_cm.png}
    \caption{Zero-shot GGUF}
    \label{fig:bloom_zero_shot_cm}
\end{subfigure}
\caption{Confusion matrices for hold-out baselines (\texttt{evaluation\_outputs/})}
\label{fig:bloom_baselines_cm}
\end{figure}

\begin{figure}[H]
\centering
\includegraphics[width=0.72\linewidth]{images/bloom_lora_cm.png}
\caption{Confusion matrix for fine-tuned Qwen LoRA on the official test split (\texttt{evaluation\_results/confusion\_matrix.png})}
\label{fig:bloom_lora_cm}
\end{figure}

\begin{figure}[H]
\centering
\includegraphics[width=0.72\linewidth]{images/bloom_svm_zero_shot_agreement.png}
\caption{SVM vs.\ zero-shot prediction agreement (\texttt{evaluation\_outputs/agreement.png})}
\label{fig:bloom_agreement}
\end{figure}

\begin{table}[H]
\centering
\caption{Bloom classification results (Qwen2.5-1.5B LoRA, full test split)}
\label{tab:bloom_results}
\renewcommand{\arraystretch}{1.2}
\begin{tabular}{|l|c|}
\hline
\textbf{Metric} & \textbf{Value} \\
\hline
Accuracy & 0.748 \\
\hline
Macro-F1 & 0.721 \\
\hline
Weighted-F1 & 0.741 \\
\hline
Mean ordinal distance & 0.452 \\
\hline
Within-one-level accuracy & 0.880 \\
\hline
Severe-error rate ($\geq$3 ordinal levels) & 0.064 \\
\hline
\end{tabular}
\end{table}

\subsection{Student RAG and Retrieval}
Table~\ref{tab:rag_results} reports the local academic QA benchmark (\texttt{results/qwen\_rag\_eval.json}).

\begin{table}[H]
\centering
\caption{Academic QA / RAG evaluation (Qwen GGUF, $n{=}10$)}
\label{tab:rag_results}
\renewcommand{\arraystretch}{1.2}
\begin{tabular}{|l|c|}
\hline
\textbf{Metric} & \textbf{Mean} \\
\hline
Exact match & 0.800 \\
\hline
Token-F1 & 0.914 \\
\hline
Retrieval hit@1 & 0.900 \\
\hline
Retrieval hit@3 & 1.000 \\
\hline
Hallucination proxy rate & 0.000 \\
\hline
\end{tabular}
\end{table}

Extractive multimodal smoke tests (\texttt{results/multimodal\_rag\_eval.json}) achieved 100\% answer accuracy on PDF and image cases when OCR dependencies are available.

\subsection{Privacy and Federated Guard}
Table~\ref{tab:privacy_results} shows PrivacyGuard metrics after \texttt{privacy/evaluate\_privacy\_guard.py}.

\begin{table}[H]
\centering
\caption{PrivacyGuard evaluation}
\label{tab:privacy_results}
\renewcommand{\arraystretch}{1.2}
\begin{tabular}{|l|c|}
\hline
\textbf{Metric} & \textbf{Value} \\
\hline
Student attack block rate & 1.000 \\
\hline
Student benign allow rate & 0.867 \\
\hline
Teacher moderation allow rate & 1.000 \\
\hline
\end{tabular}
\end{table}

All six attack categories (direct reconstruction, paraphrase probe, semantic reconstruction, etc.) achieved 100\% block rate in the evaluated taxonomy. The federated privacy guard is trained via FedAvg with clipped updates (\texttt{privacy/train\_federated\_privacy\_guard.py}).

\section{Comparison with Baseline / Existing Approaches}
On the shared hold-out split, TF-IDF + LinearSVC achieves the highest accuracy (0.839) because Bloom labels correlate with surface cues; zero-shot Qwen GGUF scores 0.441 accuracy and 1.074 mean ordinal distance. The fine-tuned LoRA classifier (0.748 accuracy, 0.721 macro-F1 on the official test set) closes most of that gap and beats zero-shot by a wide margin, showing the value of domain fine-tuning over prompt-only classification. Compared with unguarded RAG, PrivacyGuard blocks reconstruction prompts while preserving most benign study queries.

\section{Discussion on Findings}
Results show that a small on-device stack can support grounded study assistance and teacher Bloom moderation. Ordinal within-one accuracy above 88\% indicates most errors are adjacent Bloom levels---acceptable for moderation support. Perfect attack block rate on the defined suite demonstrates the value of layered guards, though this is not a proof of perfect privacy. Federated LoRA simulation enables teacher sites to contribute adapter updates without uploading raw exam text.

\section{Limitations}
\begin{itemize}
    \item Benchmarks are modest in size; broader course QA and cross-institution data are needed.
    \item Federated LoRA and privacy FL are simulated on one machine; production needs TLS and secure aggregation.
    \item OCR/multimodal tests depend on Tesseract/Pillow installation.
    \item GGUF context and RAM exceed 1\,GB host memory during inference despite sub-1\,GB weight file.
\end{itemize}

\section{Summary}
Quantitative results support the project objectives: usable Bloom classification, strong local RAG retrieval, and high adversarial block rate with reasonable benign utility. Reproducible artifacts are stored under \texttt{results/} and consolidated in \texttt{results/unified\_results\_table.csv}.
